% 张三慧第四版V2第17章《磁场和它的源》：独立候选解析
% 依据教材PDF物理页142--144的题文和图17.38--17.46。
% 尚待第三方物理审核、原图矢量描摹与辅导书逐题差异核对。
\section*{第17章\quad 磁场和它的源：独立解析}

\paragraph{习题 17.1}
采用毕奥—萨伐尔定律
\[
 d\boldsymbol B=\frac{\mu_0 I}{4\pi}
 \frac{d\boldsymbol\ell\times\boldsymbol r}{r^3}.
\]
约定图面外为 \(+\hat{\boldsymbol z}\)。
(a) 水平段上所有电流元与
\(P\) 的连线平行，故对 \(P\) 无贡献。
竖直半无限段离 \(P\) 垂距 \(a\)，
由半无限直导线
\(B=\mu_0I/(4\pi a)\)；
电流向下，\(P\) 在其右侧，
方向按右手法则为\emph{出纸面}：
\[
 \boxed{\boldsymbol B_a=
 \frac{\mu_0I}{4\pi a}\hat{\boldsymbol z}}.
\]
(b) 半圆弧半径 \(r\)，电流顺时针，
弧在中心的贡献
\(B_{\rm arc}=\mu_0I/(4r)\) 入纸面。
两条半无限平行直导线
各贡献 \(\mu_0I/(4\pi r)\)，
在 \(P\) 处同样入纸面，所以
\[
 \boxed{\boldsymbol B_b=
 -\left(\frac{\mu_0I}{4r}
 +\frac{\mu_0I}{2\pi r}\right)
 \hat{\boldsymbol z}
 =-\frac{\mu_0I(\pi+2)}{4\pi r}
 \hat{\boldsymbol z}}.
\]
(c) 等边三角形内心到边的垂距
\(h=a\sqrt3/6\)。
每边两端相对垂线的角均 \(60^\circ\)，
有限直导线公式
\(B_{\rm side}=\mu_0I(2\sin60^\circ)/(4\pi h)
 =3\mu_0I/(2\pi a)\)；
三边顺时针电流均入纸面：
\[
 \boxed{\boldsymbol B_c=
 -\frac{9\mu_0I}{2\pi a}\hat{\boldsymbol z}}.
\]
三种方向均由图17.38的箭头决定，
需保留图示版本以防方向倒置。

\paragraph{习题 17.2}
将输电线视为无限长直流导线，
地面正下方到线的距离 \(r=25\,\mathrm m\)，
\[
 \boxed{B_{\rm line}=\frac{\mu_0I}{2\pi r}
 =\frac{(4\pi\times10^{-7})(1.8\times10^3)}
 {2\pi(25)}
 =1.44\times10^{-5}\,\mathrm T}.
\]
当地地磁场
\(B_E=0.6\times10^{-4}\,\mathrm T
 =6.0\times10^{-5}\,\mathrm T\)，
\[
 \boxed{B_{\rm line}/B_E=0.240}.
\]
其大小是地磁场的约24%，但\emph{合磁场}
取决于两矢量夹角，不能简单将两大小相加。

\paragraph{习题 17.3}
距电流 \(r=0.80\,\mathrm m\)，
\(I=20\,\mathrm A\)，
长直导线产生
\[
 \boxed{B_w=\frac{\mu_0I}{2\pi r}
 =5.00\times10^{-6}\,\mathrm T}.
\]
地磁水平分量
\(B_H=0.18\times10^{-4}
 =1.80\times10^{-5}\,\mathrm T\) 指北。
导线水平磁场若与 \(B_H\) 垂直，
磁针偏角
\[
 \boxed{\tan\theta=\frac{B_w}{B_H}
 =\frac5{18}},\qquad
 \boxed{|\theta|=\arctan(5/18)\approx15.5^\circ}.
\]
为给定导线磁场\emph{大小}产生最大偏角，
其水平磁场方向应尽可能与地磁水平分量
夹成相应有利角度；若题目要求
导线与地磁垂直，这是其特定几何值。
严格优化\emph{自由方向}时，
\(\sin\theta_{\max}=B_w/B_H\)（因 \(B_w<B_H\)），
故
\(\theta_{\max}\approx16.1^\circ\)；
需区分「正交几何」与「最有利几何」两问。
\textbf{源图审核风险：}本题原文的导线在船上的
实际摆放可能约束方向，
待与辅导书核对可达到哪个最大角。
不应把随意定向的优化值当成源题直接事实。

\paragraph{习题 17.4}
设整环半径 \(R\)，
从进线点A至出线点C分成两段，
中心角分别为
\(\theta_1,\theta_2\)，
\(\theta_1+\theta_2=2\pi\)。
同种均匀导线两段电阻正比于弧长，
\(R_1/R_2=\theta_1/\theta_2\)。
两条弧并联，故电流反比分流：
\(I_1/I_2=\theta_2/\theta_1\)，
即
\(I_1\theta_1=I_2\theta_2\)。
环中心一段载流圆弧的场
\(B_{\rm arc}=\mu_0 I\theta/(4\pi R)\)，
两弧从A到C绕向相反，方向相反：
\[
 \boxed{B_O=\frac{\mu_0}{4\pi R}
 (I_1\theta_1-I_2\theta_2)=0}.
\]
径向引线对O处为零，
因为 \(d\boldsymbol\ell\parallel\boldsymbol r\)。
若环材料或截面两段不同，
该抵消结论不再一般成立。

\paragraph{习题 17.5}
两个直线电流在两线之间的磁场方向相同
（按图均出纸面）。
两线距离 \(d=0.40\,\mathrm m\)，
中点离两线 \(r=d/2=0.20\,\mathrm m\)：
\[
 \boxed{B_{\rm mid}
 =2\frac{\mu_0I}{2\pi r}
 =\frac{\mu_0I}{\pi r}
 =4.00\times10^{-5}\,\mathrm T}
\]
向纸外。
取 \(x\) 从左线到右线，
灰色矩形覆盖
\(x\in[r_1,d-r_3]=[0.10,0.30]\) m，
长 \(l=0.25\,\mathrm m\)。
其中
\[
 B(x)=\frac{\mu_0I}{2\pi}
 \left(\frac1x+\frac1{d-x}\right).
\]
故穿过矩形的正向磁通量
\[
 \begin{aligned}
 \Phi&=l\int_{r_1}^{d-r_3}B(x)\,dx\\
 &=\frac{\mu_0Il}{2\pi}
 \ln\frac{(d-r_1)(d-r_3)}{r_1r_3}.
 \end{aligned}
\]
代入 \(r_1=r_3=0.10\,\mathrm m\)
得到
\[
 \boxed{\Phi=\frac{\mu_0Il}{\pi}\ln3
 \approx2.20\times10^{-6}\,\mathrm{Wb}}
\]
（面积法向取纸外）。
跨越两线外侧时磁场方向与上式不同，
故不可用同一加号代替所有区域。

\paragraph{习题 17.6}
图17.41的半圆半径 \(R\)，
其轴取为 \(x\)，观察点
\(P=(x,0,0)\)。
令半圆位于 \(yz\) 平面，
参数化为
\(\boldsymbol r_{\rm wire}
 =(0,R\cos\varphi,R\sin\varphi)\)，
\(\varphi\in[0,\pi]\)；
电流沿 \(\varphi\) 增大方向流动
（反向电流时整个结果反号）。
令
\(D=\sqrt{x^2+R^2}\)，
则对观察点
\[
 d\boldsymbol\ell\times
 (\boldsymbol r_P-\boldsymbol r_{\rm wire})
 =(R^2,Rx\cos\varphi,Rx\sin\varphi)d\varphi.
\]
积分得到
\[
 \boxed{\boldsymbol B(P)=
 \frac{\mu_0I}{4\pi(R^2+x^2)^{3/2}}
 \big(\pi R^2\hat{\boldsymbol x}
 +2Rx\hat{\boldsymbol z}\big)}.
\]
故\emph{轴向分量}
\[
 \boxed{B_x=\frac{\mu_0IR^2}
 {4(R^2+x^2)^{3/2}}},
\]
正好是整圆环轴向场的一半；
但对于\emph{只有半圆弧}的未闭合电流路径，
一般还存在横向分量
\(\boxed{B_\perp=\mu_0IRx/
 [2\pi(R^2+x^2)^{3/2}]}\)。
源题若只要求轴向分量，可直接取 \(B_x\)；
若要求磁场矢量，不可把横向项遗漏。
实际半圆环须配接返回导线，
返回路径的贡献取决于其几何形状。

\paragraph{习题 17.7}
半径单位统一成m：
中央实心导体 \(a=0.025\)，
外筒内外半径 \(b=0.035,c=0.045\)。
进线电流 \(+I\) 匀布在内导体，
回线电流 \(-I\) 匀布在外筒。
轴向安培环路给
\(\oint\boldsymbol B\cdot d\boldsymbol\ell
 =\mu_0I_{\rm encl}(r)\)，
因而
\[
 \boxed{B(r)=\frac{\mu_0I}{2\pi}
 \begin{cases}
 r/a^2,&0\le r<a,\\
 1/r,&a\le r<b,\\
 \dfrac{1-(r^2-b^2)/(c^2-b^2)}r,
   &b\le r<c,\\
 0,&r\ge c.
 \end{cases}}
\]
磁场沿环向，
方向按进线电流右手螺旋定则。
在 \(r=a\) 处达到最大
\(B_{\max}=\mu_0I/(2\pi a)
 \approx0.040\,\mathrm T\)。
内棒中从0线性增长，间隙中按 \(1/r\)
衰减，回流筒内继续下降，
筒外在理想等流量抵消模型中为零。
绘图应在 \(a,b,c\) 各分界保持 \(B\) 连续。

\paragraph{习题 17.8}
(1) 无限长均匀实心圆柱
半径 \(R\)，总电流 \(I\)
的体电流密度为
\(J=I/(\pi R^2)\)。
对半径 \(r<R\) 的圆周，
其包围电流
\(I_{\rm enc}=J\pi r^2=Ir^2/R^2\)；
结合轴对称磁场的环路积分得
\[
 \boxed{B(r)=
 \begin{cases}
 \dfrac{\mu_0Ir}{2\pi R^2},&r\le R,\\
 \dfrac{\mu_0I}{2\pi r},&r\ge R.
 \end{cases}}
\]
(2) 无穷大平面薄电流层，
面电流密度大小为 \(j\)（A/m）。
将它分成平行的无限直电流元
\(dI=j\,dx\)，利用
\(dB=\mu_0dI/(2\pi s)\)
及在平面法线方向的投影，
对面内位置 \(x\) 积分：
\[
 B_{\rm one\ side}
 =\frac{\mu_0j}{2\pi}
 \int_{-\infty}^{\infty}
 \frac{h\,dx}{x^2+h^2}
 =\boxed{\frac{\mu_0j}{2}}
\]
（\(h>0\)，且沿平面方向
垂直于电流流向）。
在平面的另一侧大小相同、方向相反。
这个积分证明也与安培环路法相吻合。

\paragraph{习题 17.9}
取直电流沿 \(z\)，
在导线垂距 \(r\) 的点观察。
电流元位置 \(z\)，矢量距离
\(s=\sqrt{r^2+z^2}\)，
其毕奥—萨伐尔环向场
\[
 dB=\frac{\mu_0I}{4\pi}
 \frac{r\,dz}{(r^2+z^2)^{3/2}}.
\]
对无限长导线积分：
\[
 \begin{aligned}
 B&=\frac{\mu_0I}{4\pi}
 \int_{-\infty}^{\infty}
 \frac{r\,dz}{(r^2+z^2)^{3/2}}\\
 &=\frac{\mu_0I}{4\pi}
 \left[\frac{z}{r\sqrt{r^2+z^2}}\right]_{-\infty}^{+\infty}
 =\boxed{\frac{\mu_0I}{2\pi r}}.
 \end{aligned}
\]
一侧仅从0到+\(\infty\) 时则为
\(\mu_0I/(4\pi r)\)，
适用于习题17.1的半无限段。

\paragraph{习题 17.10}
设螺线管足够长、每单位轴长 \(n\) 匝，
电流 \(I\)，绕组密集而均匀。
图17.42选轴向长度为 \(a\) 的
矩形安培回路，其中一边在管内，
另一边远离管外。
忽略长直螺线管的外场，
垂直两边与局部 \(\boldsymbol B\) 垂直，
故环路积分约为
\(\oint\boldsymbol B\cdot d\boldsymbol\ell=Ba\)。
回路与螺线管缠绕电流的穿链数
\(N_{\rm link}=na\)，
\[
 Ba=\mu_0(na)I,
 \qquad \boxed{B=\mu_0nI}.
\]
有限长螺线管近端部外场不为零，
也存在端面场扩散；在这些区域
不能把无限长公式当作精确场分布。
